
\[
  K^{\bullet}(\underline{\underline{G}}_S^{n,r})=K^{\bullet}(S)[c_1,\ldots,c_n]/I_r,
\]
où \(I_r\) est un idéal homogène de \(K^{\bullet}(S)[c_1,\ldots,c_n]\), de
degré \(>r-n\). Il résulte alors de la définition (4.9.1) que \(K^{\bullet}(\underline{\underline{G}}_S^{n,\infty})\) est le séparé complété de \(K^{\bullet}(S)[c_1,\ldots,c_n]\) pour la topologie définie par les \(I_r\). Cette topologie est plus fine que la topologie définie par le poids; nous allons montrer qu'elle lui est équivalente, ce qui entrainera la proposition.

Posons \(A=K^{\bullet}(S)[c_1,\ldots,c_n]/I_r\), et soit \(C\) la \(A\)-algèbre engendrée par des générateurs \(x_1,\ldots,x_n\) soumis aux relations :
\[
  \sigma_j(x_1,\ldots,x_n)=c_j,
\]
où les \(\sigma_j\) désignent les fonctions symétriques élémentaires. D'après 4.4, l'homomorphisme canonique \(A\to C\) est injectif. On obtient :
\[
  \gamma_t(F_1-n)
  =1+c_1t+\ldots+c_nt^n
  =\prod_{i=1}^n(1+x_it).
\]
Comme \(1/\gamma_t(F_1-n)\) est un polynome, il en est de même pour \(1/(1+x_it)\) pour tout \(i\). Par suite, les \(x_i\) sont nilpotents, et il existe un entier \(N\) tel que tout polynome de degré \(>N\) par rapport aux \(x_i\), à coefficients dans \(K^{\bullet}(S)\), soit nul. Par conséquent, tout polynome de \(A\), de poids \(>N\), est nul, i.e. tout polynome de \(K^{\bullet}(S)[c_1,\ldots,c_n]\) de poids \(>N\) appartient à \(I_r\). Il s'ensuit que les deux topologies considérées sont équivalentes, ce qui achève la démonstration.


\section*{5. Applications aux fibrés projectifs; étude de \texorpdfstring{$f_*$}{f*}}

5.1.
Dans ce paragraphe, $S$ est un schéma quasi-compact, $\underline{E}$ un $\underline{O}_S$-Module localement libre de rang $r+1$, $X$ le fibré projectif associé à $\underline{E}$. Sur $X$, on a la suite exacte canonique :
\[
  0\longrightarrow \underline{F}\longrightarrow f^*(\underline{E})\longrightarrow \underline{O}_X(1)\longrightarrow 0,
\]
